\lyseis{Α5.1}{Αλγεβρική παράσταση και επιμεριστική ιδιότητα}
\sol{1} α) αριθμητική \quad β) αλγεβρική \quad γ) τρεις όρους ($2x$, $6\alpha$ και $-2$) \quad δ) $12-3^2=12-9=3$.
\sol{2} α) $2x$ \quad β) $3x-2$ \quad γ) $x+1$ \quad δ) $2x$ \quad ε) $2x+1$ \quad στ) $\dfrac{x}{2}+3$ \quad ζ) $x^2+y^2$ \quad η) $(x+y)^2$.
\sol{3} α) Το πενταπλάσιο του $\alpha$. \quad β) Ο $\alpha$ αυξημένος κατά $7$. \quad γ) Το ένα τρίτο του $\alpha$ ελαττωμένο κατά $1$. \quad δ) Το γινόμενο των $\alpha$ και $\beta$ αυξημένο κατά $1$. \quad ε) Το διπλάσιο του αθροίσματος των $\alpha$ και $\beta$.
\sol{4} α) $2\cdot5+6=16$ \quad β) $3\cdot(-2)-4=-10$ \quad γ) $4^2-3\cdot4=16-12=4$ \quad δ) $10-\dfrac{6}{2}=7$ \quad ε) $3\cdot(-2)+3=-6+3=-3$.
\sol{5} α) $(3+5-6)x=2x$ \quad β) $(7-2+1)\alpha=6\alpha$ \quad γ) $6x+4y$ \quad δ) $(9-4-5)\omega=0$ \quad ε) $2\kappa+9$.
\sol{6} α) Όροι $2x^2$, $-x$, $8$. Συντελεστής του $x^2$ είναι το $2$ και του $x$ το $-1$. \quad β) Όροι $5\alpha$, $-3\beta$, $4$. Συντελεστής του $\alpha$ είναι το $5$ και του $\beta$ το $-3$.
\sol{7} α) $3x+12$ \quad β) $10-5x$ \quad γ) $6x+12$ \quad δ) $-2x+6$ \quad ε) $8\alpha+4\beta$ \quad στ) $-3+y$.
\sol{8} α) Λάθος: $2x+6$. \quad β) Λάθος: $x+x=2x$. \quad γ) Σωστό. \quad δ) Σωστό. \quad ε) Σωστό. \quad στ) Λάθος: οι όροι έχουν διαφορετικά γράμματα και δεν ενώνονται.
\sol{9} α) $8\cdot(50+1)=400+8=408$ \quad β) $7\cdot(100-1)=700-7=693$ \quad γ) $25\cdot(100+2)=2500+50=2550$ \quad δ) $9\cdot(50-3)=450-27=423$.
\sol{10} α) Τετράγωνο $4x$, τρίγωνο $3x$, ορθογώνιο $2x+2(x+3)=2x+2x+6=4x+6$. \quad β) Για $x=1$: $4,\ 3,\ 10$. Για $x=2$: $8,\ 6,\ 14$. Για $x=5$: $20,\ 15,\ 26$. Για $x=10$: $40,\ 30,\ 46$.
\sol{11} Λογαριασμός $7+0{,}05x$ ευρώ. Για $40$ λεπτά $7+2=9$ €, για $100$ λεπτά $7+5=12$ €, για $250$ λεπτά $7+12{,}5=19{,}5$ €.
\sol{12} Ως ένα ορθογώνιο $4(x+3)$, ως δύο ορθογώνια $4x+4\cdot3=4x+12$. Για $x=5$: $4\cdot8=32$ και $20+12=32$ m$^2$.
\sol{13} Κόστος $3x+2y$. Για $x=7$ και $y=4{,}5$: $21+9=30$ €. Με έκπτωση $2$ €: $3x+2y-2=28$ €.
\sol{14} α) $2x+6+4x-4=6x+2$ \quad β) $10\alpha-5\beta-3\alpha-3\beta=7\alpha-8\beta$ \quad γ) $3x+6-x+4=2x+10$ \quad δ) $7-2x+6=13-2x$.
\sol{15} α) $A=4x+6x-15-7x+6=3x-9$. \quad β) Για $x=4$: $3\cdot4-9=3$. Στην αρχική: $16+3\cdot3-(28-6)=16+9-22=3$. Συμφωνούν.
\sol{16} α) Αδελφός $x+3$, μητέρα $(x+3)+25=x+28$. \quad β) $x+(x+3)+(x+28)=3x+31$. \quad γ) $3\cdot11+31=64$ έτη (έλεγχος: $11+14+39=64$).
\sol{17} $2(x+7)-4-2x=2x+14-4-2x=10$. Τα $x$ φεύγουν και μένει πάντα $10$, όποιον αριθμό και να σκεφτεί κανείς.
\sol{18} Σωστά: $3(x+4)=3x+12$ (το $3$ πολλαπλασιάζει και το $4$) και $-2(x-5)=-2x+10$ (το $-2$ επί $-5$ δίνει $+10$). Η $4x+3y$ δεν συντομεύεται, μένει όπως είναι, γιατί οι όροι έχουν διαφορετικά γράμματα.
\sol{19} α) $3x+3y=3(x+y)=30$ \quad β) $2(x+y)-5=20-5=15$ \quad γ) $-x-y=-(x+y)=-10$.
\sol{20} α) Το μεγάλο ορθογώνιο έχει $3\cdot6=18$ τετράγωνα. Τα δύο μικρά έχουν $3\cdot4=12$ και $3\cdot2=6$ τετράγωνα και $12+6=18$. \quad β) $3\cdot(x+2)=3x+6$.
\sol{21} Για $x=0,1,2,3$: α) $6,\ 8,\ 10,\ 12$ \quad β) $6,\ 8,\ 10,\ 12$ \quad γ) $3,\ 5,\ 7,\ 9$ \quad δ) $6,\ 8,\ 10,\ 12$. Οι α, β, δ έχουν πάντα την ίδια τιμή: $2(x+3)=2x+6$ (επιμεριστική) και $x+3+x+3=2x+6$ (αναγωγή). Η γ διαφέρει, γιατί είναι η $2x+3$.
\sol{22} Ελεύθερη απάντηση. Π.χ. «Πέντε εισιτήρια των $x$ ευρώ και $12$ € για τη συμμετοχή» ή «Πέντε ίσα κουτιά βιβλίων των $x$ κιλών και ένα δέμα των $12$ κιλών». Για $x=3$ η τιμή είναι $15+12=27$.
\sol{23} Κόστος $0{,}8\alpha+1{,}1\beta$. Για $\alpha=5$ και $\beta=3$: $4+3{,}3=7{,}3$ €.
\sol{24} Ελεύθερη απάντηση. Π.χ. «Σκέψου έναν αριθμό, πρόσθεσε $3$, διπλασίασε, αφαίρεσε $6$: βρίσκεις το διπλάσιο του αριθμού». Απόδειξη: $2(x+3)-6=2x+6-6=2x$.
